\chapter{Singular Value Decomposition}

\noindent\textit{Sources: Justin Wyss-Gallifent, \textit{Singular Value Decomposition}; Lay \& McDonald, \textit{Linear Algebra and Its Applications}, \S7.4.} The Spectral Theorem (Ch.~7.1) is used to construct the SVD.

\begin{remark}[Suggested Problems 4]
\textbf{Quiz:} Lay \S7.4.5--12; Justin 8.1. \textbf{Exam:} Lay \S7.4.13--18.
\end{remark}

%% ================================================================
\section{Introduction}
%% ================================================================

Factoring a matrix means writing it as a product of other matrices. For example, if $A$ is diagonalizable we can write $A=PDP^{-1}$, where $D$ contains eigenvalues and $P$ contains eigenvectors.

The \textbf{singular value decomposition (SVD)} is another factorization used in data compression, matrix approximation, pseudoinverses, signal analysis, and recognition.

%% ================================================================
\section{Definitions}
%% ================================================================

\begin{definition}[Singular Value Decomposition]
For any $m\times n$ real matrix $A$, the \textbf{singular value decomposition} of $A$ is a factorization
\[
A = U\Sigma V^T
\]
where:
\begin{itemize}
  \item $U$ is an $m\times m$ orthogonal matrix.
  \item $\Sigma$ is an $m\times n$ rectangular diagonal matrix with diagonal entries $\sigma_1,\ldots,\sigma_{\min(m,n)}$ all nonnegative and in nonincreasing order.
  \item $V$ is an $n\times n$ orthogonal matrix.
\end{itemize}
\end{definition}

\begin{definition}[Orthogonal Matrix]
An $n\times n$ matrix $U$ is \textbf{orthogonal} if its columns are unit vectors all perpendicular to one another. Equivalently: columns form an orthonormal basis of $\mathbb{R}^n$, and $U^TU=UU^T=I$.
\end{definition}

\begin{definition}[Rectangular Diagonal Matrix]
An $m\times n$ \textbf{rectangular diagonal} matrix has nonzero entries only in positions $(1,1)$, $(2,2)$, etc. Extra rows or columns (if $m\neq n$) are all zeros.
\end{definition}

%% ================================================================
\section{Constructing the SVD}
%% ================================================================

\subsection{Key Observations}

\textbf{Observation 1.} If $k\leq\min(m,n)$ singular values are nonzero, then $A=U\Sigma V^T$ depends only on the first $k$ columns of $U$ and $V$. Moreover:
\[
A\bar{v}_i = \sigma_i\bar{u}_i,\qquad A^T\bar{u}_i = \sigma_i\bar{v}_i\quad\text{for }1\leq i\leq k.
\]

\textbf{Observation 2.} Consider $AA^T$ (size $m\times m$) and $A^TA$ (size $n\times n$):
\begin{itemize}
  \item Both are symmetric; by the Spectral Theorem they have orthogonal eigenvectors.
  \item All eigenvalues are nonnegative: $\|A\bar{v}\|^2 = \bar{v}^T A^TA\bar{v} = \lambda\|\bar{v}\|^2 = \lambda$.
  \item There is a bijection between eigenpairs of $AA^T$ and $A^TA$ for nonzero eigenvalues:
  \[
  (\sigma^2,\bar{u})\mapsto \bigl(\sigma^2,\tfrac{1}{\sigma}A^T\bar{u}\bigr),\qquad
  (\sigma^2,\bar{v})\mapsto \bigl(\sigma^2,\tfrac{1}{\sigma}A\bar{v}\bigr).
  \]
\end{itemize}

\textbf{Observation 3.} For each positive eigenvalue $\sigma_i^2$ common to both, we have paired unit eigenvectors $\bar{u}_i$ (left singular vector) and $\bar{v}_i$ (right singular vector) satisfying $A\bar{v}_i=\sigma_i\bar{u}_i$ and $A^T\bar{u}_i=\sigma_i\bar{v}_i$.

\subsection{Construction in Brief}

\begin{enumerate}
  \item Find the positive eigenvalues of $AA^T$ (or $A^TA$). Denote them $\sigma_1^2\geq\sigma_2^2\geq\cdots\geq\sigma_k^2$.
  \item For each $\sigma_i^2$, find unit eigenvector $\bar{u}_i$ of $AA^T$ and set $\bar{v}_i = \frac{1}{\sigma_i}A^T\bar{u}_i$ (or vice versa). Handle multiplicity by spanning eigenspaces.
  \item Build $\Sigma$ with $\sigma_i$ on the diagonal; build $U$ from $\bar{u}_i$ and $V$ from $\bar{v}_i$. Fill remaining columns of $U$ and $V$ with eigenvectors for eigenvalue $0$ to complete orthonormal bases.
\end{enumerate}

\begin{remark}[Which matrix should you use?]
If $A$ is $m\times n$, then $A^TA$ is $n\times n$ and $AA^T$ is $m\times m$.
\begin{itemize}
  \item Use the smaller one first when doing calculations by hand.
  \item The eigenvectors of $A^TA$ give the \textbf{right} singular vectors ($V$).
  \item The eigenvectors of $AA^T$ give the \textbf{left} singular vectors ($U$).
\end{itemize}
\end{remark}

%% ================================================================
\section{What to Show on an SVD Problem}
%% ================================================================

When the problem asks you to \emph{find an SVD}, the expected structure is:
\begin{enumerate}
  \item Compute $A^TA$ or $AA^T$.
  \item Find its eigenvalues and convert them into singular values $\sigma_i=\sqrt{\lambda_i}$.
  \item Find orthonormal eigenvectors and assemble $V$ (or $U$).
  \item Use $u_i=\frac{1}{\sigma_i}Av_i$ or $v_i=\frac{1}{\sigma_i}A^Tu_i$.
  \item Write the final factorization $A=U\Sigma V^T$ and, if requested, verify by multiplication.
\end{enumerate}

\begin{remark}
If a singular value is $0$, the formula $u_i=\frac{1}{\sigma_i}Av_i$ does not apply. In that case, complete the relevant orthonormal basis using eigenvectors corresponding to the eigenvalue $0$.
\end{remark}

%% ================================================================
\section{Worked Example}
%% ================================================================

Find an SVD of
\[
A=\begin{pmatrix}
-2 & 0\\
0 & 0
\end{pmatrix}.
\]

\textbf{Step 1: Compute $A^TA$.}
\[
A^TA=
\begin{pmatrix}
4 & 0\\
0 & 0
\end{pmatrix}.
\]
So the eigenvalues are $\lambda_1=4$ and $\lambda_2=0$, hence the singular values are
\[
\sigma_1=2,\qquad \sigma_2=0.
\]

\textbf{Step 2: Find $V$.}
An orthonormal eigenbasis for $A^TA$ is
\[
v_1=\begin{pmatrix}1\\0\end{pmatrix},
\qquad
v_2=\begin{pmatrix}0\\1\end{pmatrix},
\]
so we may take
\[
V=\begin{pmatrix}
1 & 0\\
0 & 1
\end{pmatrix}.
\]

\textbf{Step 3: Find the nonzero left singular vector.}
\[
u_1=\frac{1}{\sigma_1}Av_1
=\frac{1}{2}\begin{pmatrix}-2\\0\end{pmatrix}
=\begin{pmatrix}-1\\0\end{pmatrix}.
\]
To complete an orthonormal basis of $\mathbb{R}^2$, choose
\[
u_2=\begin{pmatrix}0\\1\end{pmatrix}.
\]
Thus
\[
U=\begin{pmatrix}
-1 & 0\\
0 & 1
\end{pmatrix}.
\]

\textbf{Step 4: Build $\Sigma$.}
\[
\Sigma=\begin{pmatrix}
2 & 0\\
0 & 0
\end{pmatrix}.
\]

Therefore one SVD is
\[
\boxed{
A=U\Sigma V^T
=
\begin{pmatrix}
-1 & 0\\
0 & 1
\end{pmatrix}
\begin{pmatrix}
2 & 0\\
0 & 0
\end{pmatrix}
\begin{pmatrix}
1 & 0\\
0 & 1
\end{pmatrix}.
}
\]

%% ================================================================
\section{Best Low-Rank Approximation}
%% ================================================================

\begin{theorem}[Eckart--Young]
Among all matrices of rank $\leq k$, the best approximation to $A$ (in Frobenius norm) that preserves the most variance is
\[
A_k = \sum_{i=1}^k \sigma_i\bar{u}_i\bar{v}_i^T = U\Sigma_k V^T,
\]
where $\Sigma_k$ is $\Sigma$ with $\sigma_{k+1},\ldots,\sigma_{\min(m,n)}$ set to zero.
\end{theorem}

The proportion of variance preserved is
\[
\frac{\sigma_1^2+\cdots+\sigma_k^2}{\sigma_1^2+\cdots+\sigma_{\min(m,n)}^2}.
\]

\begin{remark}
This same quantity is what appears in image-compression problems as the \textbf{quality} or \textbf{variance preserved}. Keeping the first $k$ singular values gives the best rank-$k$ approximation in this sense.
\end{remark}

%% ================================================================
\section{Fundamental Subspaces and Reduced SVD (Lay-McDonald \S7.4)}
%% ================================================================

Given $A=U\Sigma V^T$ with rank $r$: $\{u_1,\ldots,u_r\}$ is an orthonormal basis for $\operatorname{Col} A$; $\{u_{r+1},\ldots,u_m\}$ for $\operatorname{Nul} A^T$; $\{v_1,\ldots,v_r\}$ for $\operatorname{Row} A$; $\{v_{r+1},\ldots,v_n\}$ for $\operatorname{Nul} A$.

\textbf{Reduced SVD:} $A=U_r D V_r^T$ where $U_r$, $V_r$ contain only the first $r$ columns. The \textbf{pseudoinverse} is $A^+=V_r D^{-1} U_r^T$; $\hat{x}=A^+ b$ is the least-squares solution of $Ax=b$ with smallest norm.

\textbf{Condition number:} For invertible $A$, $\kappa(A)=\sigma_1/\sigma_n$ measures sensitivity of $Ax=b$ to perturbations.
